\chapter{Additional Mathematical Derivations}

In this appendix, we consider mathematical derivations of results used throughout the thesis. These were derivations that were not relevant to the main thrust of considering the violation of conditional site-independence in models, and thus were omitted from the main text and included here.


\section{Stationary Distribution of CTMC}
\label{app:stationaryprob}
Here we prove that the binary CTMC found in Figure \ref{fig:markovchain} has stationary distribution \begin{equation}
\label{eq:stationarydistributionpi}
\boldsymbol{\pi} = (\pi_0, \pi_1) = \left(\frac{q_{10}}{q_{01}+q_{10}}, \frac{q_{01}}{q_{01}+q_{10}}\right).
\end{equation}

We start by considering $\mathbf{G}$. We note that any stationary distribution $\boldsymbol{\pi}$ satisfies $\boldsymbol{\pi}\mathbf{G} = \boldsymbol{0}$. Clearly, this is true of (\ref{eq:stationarydistributionpi}):
\begin{align}
    \boldsymbol{\pi}\mathbf{G} &= \begin{pmatrix}\frac{q_{10}}{q_{01}+q_{10}} & \frac{q_{01}}{q_{01}+q_{10}}\end{pmatrix} \begin{pmatrix}
        -q_{01} & q_{01} \\ q_{10} & -q_{10}
    \end{pmatrix} \\
    &= \begin{pmatrix}\frac{-q_{10}q_{01} +  q_{01}q_{10}}{q_{01}+q_{10}} & \frac{q_{10}q_{01} - q_{01}q_{10}}{q_{01}q_{10}}\end{pmatrix} = \mathbf{0}.
\end{align}

Thus $\boldsymbol{\pi}$ is a stationary distribution of the CTMC with generator $\mathbf{G}$. We now consider the uniqueness of $\boldsymbol{\pi}$. Clearly both states are connected, making the Markov chain generated from the CTMC irreducible. Since the state space is finite, we have that the CTMC must admit a unique stationary distribution \cite{Douc2019}. This justification is also used for Proposition \ref{prop:ctmcprod} to show that $\tilde{\boldsymbol{\pi}}$ is also stationary.

\section{Markov Property for Stochastic Process $\mathbf{D}_{l, k}$}
\label{app:ctmcprodmarkov}

We now turn to consider the fact that $\mathbf{D}_{l, k}$, the resulting binary encoding of the element-wise product of dependency $\mathbf{D}_{l-1, j}$ and simulated CTMC $\tilde{\mathbf{D}}_{l, k}$ as considered by (\ref{eq:dependencystructure}) in Section \ref{sec:ctmcproduct}, is not in fact a CTMC, because it does not satisfy the Markov property.

We define a partition on state space $\mathcal{S}$:
\begin{equation}
\label{eq:lumpability}
    \mathcal{P} = \{\{00, 01, 10\},\{11\}\}
\end{equation}
This is a clear case of reducing the size of the state space. We employ the idea of \textit{Markov lumpability} \cite{Kemedy1976}:
\begin{definition} [Lumpability]
Suppose $(X_t)_{t\geq 0}$ is a CTMC with finite state space $\mathcal{S}$ and generator matrix $\mathbf{G} = (g_{xy})_{x, y \in \mathcal{S}}$, and $\mathcal{P} := \{P_1, ..., P_m\}$ is a partition of $\mathcal{S}$, then aggregated process $Y_t = i$ where $X_t \in B_i$, then we say that the CTMC is lumpable with respect to partition $\mathcal{P}$ if for every pair of blocks $P_i, P_j$, $\forall x, x' \in P_i$:
\begin{equation}
    \sum_{y\in P_j}g_{xy} = \sum_{y \in P_j}q_{x'y}.
\end{equation}
\end{definition}
This condition can be interpreted as ensuring that the jump probability from block $P_i$ to block $P_j$ is constant, regardless of the state within $P_i$. 

Clearly, this is an important property the Markov chain must have in order for the reduced state space stochastic process to be considered a Markov chain as well. We state the formal result without proof (of which can be found in \cite{Kemedy1976}):
\begin{proposition} [Resulting process of lumped Markov chain is Markovian]
\label{prop:markovprocesses}
Suppose $(X_t)_{t\geq 0}$ is a CTMC with finite state space $\mathcal{S}$ and generator matrix $\mathbf{G} = (g_{xy})_{x, y \in \mathcal{S}}$, and $\mathcal{P} := \{P_1, ..., P_m\}$ is a partition of $\mathcal{S}$. Then the resulting aggregated process $(Y_t)_{t\geq0}$ is also a CTMC if and only if the CTMC is lumpable.
\end{proposition}

As such, in order to prove that $\mathbf{D}_{l, k}$ does not satisfy the Markov property, we need only show that $\mathbf{B}_{l, k}$ is not lumpable with respect to the partition (\ref{eq:lumpability}). This is true, as the generator $\mathbf{G}_{\text{prod}}$ as defined in (\ref{eq:ctmcprodgenerator}) shows: we note that both $\{01, 10\}$ have an equivalent jump probability to $\{11\}$, but that is not the case with state $\{00\}$, which has a zero transition probability to $\{11\}$. As such, from Proposition \ref{prop:markovprocesses}, we can conclude that the lumped chain $\mathbf{D}_{l, k}$ with respect to the partition $\mathcal{P}$ defined in \ref{eq:lumpability} indeed does not follow the Markov property.